\scanpage{20}\mspage{16}
% @ top, lines 1–4 | text | donc t = {u,v,w}
donc $t=\{u,v,w\}$ est unique de son espèce. Si $x$ et $y$ sont non liés, on sait qu'il existe un \inserted{(et $=$ deux)} $a\in E(u)^*$ tel que $a$ lié à la fois à $x$ et $y$, donc $t$ non unique de son espèce, cqfd.

% @ upper middle, corollary | text | Corollaire Le groupe
\uline{Corollaire} Le groupe \deleted{\unsure{des}} $W_u$ \inserted{$\subset W$} des autom de $C$ qui fixent $\unsure{a|u}$ est \deleted{\unsure{un}} isom au groupe des couples $(\unsure{\varphi|\rho},\psi)$ d'un autom $\varphi\in\unsure{\mathfrak S|\mathscr S}_{Q_u}$ de \unsure{$P_u$|$Q_u$}, et d'un $\overline{\varphi}$ autom. du \deleted{\unsure{$P$}} $\Psi_{Q_u}$ torseur $P$, où $\overline{\varphi}$ est l'autom de $\Psi_{Q_u}$ induit par $\varphi$. Donc on a une suite exacte (splittée)

% @ centre, exact sequence and vertical identifications | diagram | 1, (Ψ), W, S, 1 ; F₂⁴ and S₅
\diagram{\begin{tikzcd}[ampersand replacement=\&,column sep=large,row sep=small]
1\arrow[r,dash] \& (\Psi_{Q_u})^{\deleted{\unsure{\text{torseur}}}}\arrow[r] \& W_u\arrow[r,dash] \& \mathfrak S_{Q_u}\arrow[r,dash] \&1\\
\&\mathbb F_2^4\arrow[u,phantom,"\simeq"{description,rotate=90}] \& \&\mathfrak S_5\arrow[u,phantom,"\simeq"{description,rotate=90}] \&
\end{tikzcd}}

% @ just below fold | text | donc Wᵤ ≃ groupe de Weyl
donc $W_u\simeq$ groupe de Weyl $W_{D_5}$ de $D_5$ $(\simeq\unsure{\mathrm{so}(10)|\mathrm{SO}(10)})$ qui apparait donc comme ss-groupe d'indice 27 dans $W$ $(\simeq W_{E_6})$.

% @ lower third, NB | text | Il faudrait expliciter
\uline{NB} Il faudrait expliciter comment la donnée de $(C,u)$ permet de définir un syst. de racines $D_5$ (plongé dans le syst. de racines $E_6$ associé à $C$) -- donc définir un ens de \inserted{40} ``bases'' (=``racines'') de $C$, stable par passage à bases opposées. Or on trouve qu'il y a $40=(10.8):2$ \inserted{$=\binom{10}{2}-5$} \deleted{couples} \inserted{paires} d'éléments de $E(u)$ non liés. Or on a le

% @ foot, lemma | text | Lemme Pour tte paire
\uline{Lemme} Pour tte paire \inserted{$\{\unsure{\zeta_1|\xi_1},\unsure{\zeta_2|\xi_2}\}$} d'éléments non liés de
% archive sheet number: 19
